\begin{abstract}
Random file splits can place near-duplicate PHP webshells in both the training and test sets, making it
difficult to assess performance on unseen sources. We reproduce GlareShell, a static detector that classifies
PHP files using interprocedural control-flow graphs, and evaluate it on a reconstructed benchmark.
Across five seeds, the graph attention network achieves a mean F1 score of
\result{gat_paper_f1} under random splitting and \result{gat_clean_f1} when each near-duplicate group stays
in one partition. Logistic regression on TF-IDF features achieves \result{lr_clean_f1} on the grouped
splits. In the test on held-out sources, with \result{cross_source_test_benign} benign files and
\result{cross_source_test_webshell} files labelled as webshells by source, the graph model yields a false-positive rate of
\result{gat_cross_fpr_percent} and a recall of \result{gat_cross_recall}, compared with
\result{lr_cross_fpr_percent} and \result{lr_cross_recall} for the baseline. These findings support evaluating
webshell detectors on unseen projects and collections before drawing conclusions about deployment.
\end{abstract}

\begin{IEEEkeywords}
PHP webshells, graph neural networks, static analysis, reproducibility, duplicate control, source shift
\end{IEEEkeywords}
